
\documentclass{article}
\usepackage{colortbl}
\usepackage{makecell}
\usepackage{multirow}
\usepackage{supertabular}

\begin{document}

\newcounter{utterance}

\centering \large Interaction Transcript for game `hot\_air\_balloon', experiment `air\_balloon\_survival\_en\_reasoning off\_hard', episode 5 with claude{-}sonnet{-}4{-}20250514{-}no{-}reasoning{-}t1.0.
\vspace{24pt}

{ \footnotesize  \setcounter{utterance}{1}
\setlength{\tabcolsep}{0pt}
\begin{supertabular}{c@{$\;$}|p{.15\linewidth}@{}p{.15\linewidth}p{.15\linewidth}p{.15\linewidth}p{.15\linewidth}p{.15\linewidth}}
    \# & $\;$A & \multicolumn{4}{c}{Game Master} & $\;\:$B\\
    \hline

    \theutterance \stepcounter{utterance}  
    & & \multicolumn{4}{p{0.6\linewidth}}{
        \cellcolor[rgb]{0.9,0.9,0.9}{
            \makecell[{{p{\linewidth}}}]{
                \texttt{\tiny{[P1$\langle$GM]}}
                \texttt{Stai partecipando a un gioco di negoziazione collaborativa.} \\
\\ 
\texttt{Insieme a un altro partecipante, dovete accordarvi su un unico set di oggetti da conservare. Ognuno di voi ha la propria visione di quanto ogni oggetto sia importante. Non sai come l'altro partecipante valuta gli oggetti. Inoltre, ti viene fornito lo sforzo che ogni oggetto richiede.} \\
\texttt{Potete accordarvi su un set solo se lo sforzo totale degli oggetti selezionati non supera un limite condiviso:} \\
\\ 
\texttt{LIMITE = 2307} \\
\\ 
\texttt{Ecco i valori di sforzo dei singoli oggetti:} \\
\\ 
\texttt{Sforzo oggetto = \{"A21": 147, "B93": 698, "A11": 588, "B33": 43, "C08": 356, "B39": 80, "B96": 95, "C85": 743, "A82": 106, "B31": 308, "B20": 325, "A60": 255, "B09": 276, "C81": 543, "B21": 51\}} \\
\\ 
\texttt{Ecco la tua visione personale sull'importanza di ogni oggetto:} \\
\\ 
\texttt{Valori di importanza dell'oggetto = \{"A21": 138, "B93": 583, "A11": 65, "B33": 262, "C08": 121, "B39": 508, "B96": 461, "C85": 484, "A82": 668, "B31": 389, "B20": 215, "A60": 97, "B09": 500, "C81": 30, "B21": 400\}} \\
\\ 
\texttt{Obiettivo:} \\
\\ 
\texttt{Il tuo obiettivo è negoziare un set condiviso di oggetti che ti avvantaggi il più possibile (cioè, massimizzi l'importanza totale per TE), rimanendo entro il LIMITEE. Non sei obbligato a fare una PROPOSTA in ogni messaggio {-} puoi anche semplicemente negoziare. Tutte le tattiche sono permesse!} \\
\\ 
\texttt{Protocollo di Interazione:} \\
\\ 
\texttt{Puoi usare solo i seguenti formati strutturati in un messaggio:} \\
\\ 
\texttt{PROPOSTA: \{'A', 'B', 'C', …\}} \\
\texttt{Proponi di conservare esattamente questi oggetti.} \\
\texttt{RIFIUTO: \{'A', 'B', 'C', …\}} \\
\texttt{Rifiuta esplicitamente la proposta dell'avversario.} \\
\texttt{ARGOMENTAZIONE: \{'...'\}} \\
\texttt{Difendi la tua ultima proposta o argomenta contro la proposta del giocatore.} \\
\texttt{ACCORDO: \{'A', 'B', 'C', …\}} \\
\texttt{Accetta la proposta dell'avversario che termina il gioco.} \\
\\ 
\\ 
\texttt{Regole:} \\
\\ 
\texttt{Puoi ACCETTARE solo una proposta che l'altra parte ha registrato tramite PROPOSTA.} \\
\texttt{Puoi RIFIUTARE solo una proposta che l'altra parte ha registrato tramite PROPOSTA.} \\
\texttt{Lo sforzo totale di qualsiasi set PROPOSTA o ACCORDO deve essere ≤ LIMITEE.} \\
\texttt{NON rivelare i tuoi punteggi di importanza nascosti.} \\
\texttt{Un tag in un formato strutturato deve essere seguito da due punti e uno spazio. L'argomento deve essere un set python contenente 0 o più stringhe.} \\
\texttt{Quindi, deve essere della forma TAG: \{...\}} \\
\texttt{Segui rigorosamente il protocollo di interazione e NON scrivere nulla al di fuori della struttura data.} \\
\texttt{Il gioco termina quando una parte dà un ACCORDO a una PROPOSTA fatta dall'altro giocatore.} \\
\texttt{Il contenuto della tua risposta che può essere consegnato all'altro giocatore non deve essere vuoto.} \\
\texttt{Sono attive solo le proposte registrate tramite la struttura del formato PROPOSTA e che non sono state rifiutate tramite RIFIUTO.} \\
\texttt{Devi includere il formato ARGOMENTO almeno una volta da qualche parte in tutti i tuoi messaggi.} \\
            }
        }
    }
    & \\ \\

    \theutterance \stepcounter{utterance}  
    & \multicolumn{4}{p{0.6\linewidth}}{
        \cellcolor[rgb]{0.9,0.9,0.9}{
            \makecell[{{p{\linewidth}}}]{
                \texttt{\tiny{[P1$\rangle$GM]}}
                \texttt{PROPOSTA: \{'A82', 'B39', 'B96', 'B09', 'B21', 'B31', 'A21', 'B33'\}} \\
\\ 
\texttt{ARGOMENTAZIONE: \{'Ho selezionato oggetti che offrono un buon equilibrio tra valore e sforzo richiesto. Questa combinazione include alcuni oggetti ad alto valore come A82 e B39, insieme ad altri che hanno un rapporto costo{-}beneficio interessante. Il totale dello sforzo è di 1016, lasciando ancora margine per eventuali aggiustamenti se necessari.'\}} \\
            }
        }
    }
    & & \\ \\

    \theutterance \stepcounter{utterance}  
    & & \multicolumn{4}{p{0.6\linewidth}}{
        \cellcolor[rgb]{0.9,0.9,0.9}{
            \makecell[{{p{\linewidth}}}]{
                \texttt{\tiny{[GM$\rangle$P2]}}
                \texttt{Stai partecipando a un gioco di negoziazione collaborativa.} \\
\\ 
\texttt{Insieme a un altro partecipante, dovete accordarvi su un unico set di oggetti da conservare. Ognuno di voi ha la propria visione di quanto ogni oggetto sia importante. Non sai come l'altro partecipante valuta gli oggetti. Inoltre, ti viene fornito lo sforzo che ogni oggetto richiede.} \\
\texttt{Potete accordarvi su un set solo se lo sforzo totale degli oggetti selezionati non supera un limite condiviso:} \\
\\ 
\texttt{LIMITE = 2307} \\
\\ 
\texttt{Ecco i valori di sforzo dei singoli oggetti:} \\
\\ 
\texttt{Sforzo oggetto = \{"A21": 147, "B93": 698, "A11": 588, "B33": 43, "C08": 356, "B39": 80, "B96": 95, "C85": 743, "A82": 106, "B31": 308, "B20": 325, "A60": 255, "B09": 276, "C81": 543, "B21": 51\}} \\
\\ 
\texttt{Ecco la tua visione personale sull'importanza di ogni oggetto:} \\
\\ 
\texttt{Valori di importanza dell'oggetto = \{"A82": 30, "B93": 65, "B39": 97, "B09": 121, "C85": 138, "B96": 215, "B21": 262, "B31": 389, "B33": 400, "B20": 461, "A21": 484, "C08": 500, "A60": 508, "A11": 583, "C81": 668\}} \\
\\ 
\texttt{Obiettivo:} \\
\\ 
\texttt{Il tuo obiettivo è negoziare un set condiviso di oggetti che ti avvantaggi il più possibile (cioè, massimizzi l'importanza totale per TE), rimanendo entro il LIMITEE. Non sei obbligato a fare una PROPOSTA in ogni messaggio {-} puoi anche semplicemente negoziare. Tutte le tattiche sono permesse!} \\
\\ 
\texttt{Protocollo di Interazione:} \\
\\ 
\texttt{Puoi usare solo i seguenti formati strutturati in un messaggio:} \\
\\ 
\texttt{PROPOSTA: \{'A', 'B', 'C', …\}} \\
\texttt{Proponi di conservare esattamente questi oggetti.} \\
\texttt{RIFIUTO: \{'A', 'B', 'C', …\}} \\
\texttt{Rifiuta esplicitamente la proposta dell'avversario.} \\
\texttt{ARGOMENTAZIONE: \{'...'\}} \\
\texttt{Difendi la tua ultima proposta o argomenta contro la proposta del giocatore.} \\
\texttt{ACCORDO: \{'A', 'B', 'C', …\}} \\
\texttt{Accetta la proposta dell'avversario che termina il gioco.} \\
\\ 
\\ 
\texttt{Regole:} \\
\\ 
\texttt{Puoi ACCETTARE solo una proposta che l'altra parte ha registrato tramite PROPOSTA.} \\
\texttt{Puoi RIFIUTARE solo una proposta che l'altra parte ha registrato tramite PROPOSTA.} \\
\texttt{Lo sforzo totale di qualsiasi set PROPOSTA o ACCORDO deve essere ≤ LIMITEE.} \\
\texttt{NON rivelare i tuoi punteggi di importanza nascosti.} \\
\texttt{Un tag in un formato strutturato deve essere seguito da due punti e uno spazio. L'argomento deve essere un set python contenente 0 o più stringhe.} \\
\texttt{Quindi, deve essere della forma TAG: \{...\}} \\
\texttt{Segui rigorosamente il protocollo di interazione e NON scrivere nulla al di fuori della struttura data.} \\
\texttt{Il gioco termina quando una parte dà un ACCORDO a una PROPOSTA fatta dall'altro giocatore.} \\
\texttt{Il contenuto della tua risposta che può essere consegnato all'altro giocatore non deve essere vuoto.} \\
\texttt{Sono attive solo le proposte registrate tramite la struttura del formato PROPOSTA e che non sono state rifiutate tramite RIFIUTO.} \\
\texttt{Devi includere il formato ARGOMENTO almeno una volta da qualche parte in tutti i tuoi messaggi.} \\
\\ 
\\ 
\texttt{Ora riceverai il primo messaggio dell'altro giocatore.} \\
\\ 
\texttt{PROPOSTA: \{'A82', 'B39', 'B96', 'B09', 'B21', 'B31', 'A21', 'B33'\}} \\
\\ 
\texttt{ARGOMENTAZIONE: \{'Ho selezionato oggetti che offrono un buon equilibrio tra valore e sforzo richiesto. Questa combinazione include alcuni oggetti ad alto valore come A82 e B39, insieme ad altri che hanno un rapporto costo{-}beneficio interessante. Il totale dello sforzo è di 1016, lasciando ancora margine per eventuali aggiustamenti se necessari.'\}} \\
            }
        }
    }
    & \\ \\

    \theutterance \stepcounter{utterance}  
    & & & \multicolumn{4}{p{0.6\linewidth}}{
        \cellcolor[rgb]{0.9,0.9,0.9}{
            \makecell[{{p{\linewidth}}}]{
                \texttt{\tiny{[GM$\langle$P2]}}
                \texttt{RIFIUTO: \{'A82', 'B39', 'B96', 'B09', 'B21', 'B31', 'A21', 'B33'\}} \\
\\ 
\texttt{ARGOMENTAZIONE: \{'La tua proposta include diversi oggetti con rapporto valore{-}sforzo non ottimale. Possiamo utilizzare meglio il nostro budget di sforzo includendo oggetti più strategici che massimizzino il valore complessivo per entrambi.'\}} \\
\\ 
\texttt{PROPOSTA: \{'A11', 'C81', 'A60', 'C08', 'B20', 'B33', 'B21', 'B39', 'B96'\}} \\
            }
        }
    }
     \\ \\

    \theutterance \stepcounter{utterance}  
    & & \multicolumn{4}{p{0.6\linewidth}}{
        \cellcolor[rgb]{0.9,0.9,0.9}{
            \makecell[{{p{\linewidth}}}]{
                \texttt{\tiny{[P1$\langle$GM]}}
                \texttt{RIFIUTO: \{'A82', 'B39', 'B96', 'B09', 'B21', 'B31', 'A21', 'B33'\}} \\
\\ 
\texttt{ARGOMENTAZIONE: \{'La tua proposta include diversi oggetti con rapporto valore{-}sforzo non ottimale. Possiamo utilizzare meglio il nostro budget di sforzo includendo oggetti più strategici che massimizzino il valore complessivo per entrambi.'\}} \\
\\ 
\texttt{PROPOSTA: \{'A11', 'C81', 'A60', 'C08', 'B20', 'B33', 'B21', 'B39', 'B96'\}} \\
            }
        }
    }
    & \\ \\

    \theutterance \stepcounter{utterance}  
    & \multicolumn{4}{p{0.6\linewidth}}{
        \cellcolor[rgb]{0.9,0.9,0.9}{
            \makecell[{{p{\linewidth}}}]{
                \texttt{\tiny{[P1$\rangle$GM]}}
                \texttt{RIFIUTO: \{'A11', 'C81', 'A60', 'C08', 'B20', 'B33', 'B21', 'B39', 'B96'\}} \\
\\ 
\texttt{ARGOMENTAZIONE: \{'La tua proposta include oggetti come C81 e A11 che potrebbero non essere ottimali dal punto di vista del rapporto valore{-}sforzo. Inoltre, A60 richiede molto sforzo per il valore che offre. Propongo una combinazione più efficiente che massimizza il valore totale mantenendo un uso ragionevole del budget di sforzo.'\}} \\
\\ 
\texttt{PROPOSTA: \{'A82', 'B93', 'B39', 'B96', 'B09', 'B21', 'A21', 'B33'\}} \\
            }
        }
    }
    & & \\ \\

    \theutterance \stepcounter{utterance}  
    & & \multicolumn{4}{p{0.6\linewidth}}{
        \cellcolor[rgb]{0.9,0.9,0.9}{
            \makecell[{{p{\linewidth}}}]{
                \texttt{\tiny{[GM$\rangle$P2]}}
                \texttt{RIFIUTO: \{'A11', 'C81', 'A60', 'C08', 'B20', 'B33', 'B21', 'B39', 'B96'\}} \\
\\ 
\texttt{ARGOMENTAZIONE: \{'La tua proposta include oggetti come C81 e A11 che potrebbero non essere ottimali dal punto di vista del rapporto valore{-}sforzo. Inoltre, A60 richiede molto sforzo per il valore che offre. Propongo una combinazione più efficiente che massimizza il valore totale mantenendo un uso ragionevole del budget di sforzo.'\}} \\
\\ 
\texttt{PROPOSTA: \{'A82', 'B93', 'B39', 'B96', 'B09', 'B21', 'A21', 'B33'\}} \\
            }
        }
    }
    & \\ \\

    \theutterance \stepcounter{utterance}  
    & & & \multicolumn{4}{p{0.6\linewidth}}{
        \cellcolor[rgb]{0.9,0.9,0.9}{
            \makecell[{{p{\linewidth}}}]{
                \texttt{\tiny{[GM$\langle$P2]}}
                \texttt{RIFIUTO: \{'A82', 'B93', 'B39', 'B96', 'B09', 'B21', 'A21', 'B33'\}} \\
\\ 
\texttt{ARGOMENTAZIONE: \{'La tua proposta include B93 che richiede ben 698 di sforzo, lasciando poco spazio per altri oggetti di valore. Inoltre A82 ha un rapporto valore{-}sforzo molto basso. Possiamo ottenere risultati migliori con una selezione più bilanciata che non sprechi così tanto budget su singoli oggetti.'\}} \\
\\ 
\texttt{PROPOSTA: \{'A11', 'A60', 'C08', 'B20', 'B33', 'B21', 'B39', 'B96', 'B09'\}} \\
            }
        }
    }
     \\ \\

    \theutterance \stepcounter{utterance}  
    & & \multicolumn{4}{p{0.6\linewidth}}{
        \cellcolor[rgb]{0.9,0.9,0.9}{
            \makecell[{{p{\linewidth}}}]{
                \texttt{\tiny{[P1$\langle$GM]}}
                \texttt{RIFIUTO: \{'A82', 'B93', 'B39', 'B96', 'B09', 'B21', 'A21', 'B33'\}} \\
\\ 
\texttt{ARGOMENTAZIONE: \{'La tua proposta include B93 che richiede ben 698 di sforzo, lasciando poco spazio per altri oggetti di valore. Inoltre A82 ha un rapporto valore{-}sforzo molto basso. Possiamo ottenere risultati migliori con una selezione più bilanciata che non sprechi così tanto budget su singoli oggetti.'\}} \\
\\ 
\texttt{PROPOSTA: \{'A11', 'A60', 'C08', 'B20', 'B33', 'B21', 'B39', 'B96', 'B09'\}} \\
            }
        }
    }
    & \\ \\

    \theutterance \stepcounter{utterance}  
    & \multicolumn{4}{p{0.6\linewidth}}{
        \cellcolor[rgb]{0.9,0.9,0.9}{
            \makecell[{{p{\linewidth}}}]{
                \texttt{\tiny{[P1$\rangle$GM]}}
                \texttt{RIFIUTO: \{'A11', 'A60', 'C08', 'B20', 'B33', 'B21', 'B39', 'B96', 'B09'\}} \\
\\ 
\texttt{ARGOMENTAZIONE: \{'La tua proposta continua a includere A11 e A60 che hanno un rapporto valore{-}sforzo problematico. A11 richiede 588 di sforzo mentre A60 ne richiede 255, consumando una porzione significativa del budget. Dovremmo concentrarci su oggetti che offrono maggiore efficienza nella distribuzione delle risorse disponibili.'\}} \\
\\ 
\texttt{PROPOSTA: \{'A82', 'B39', 'B96', 'B09', 'B21', 'B31', 'A21', 'B33', 'C08', 'B20'\}} \\
            }
        }
    }
    & & \\ \\

    \theutterance \stepcounter{utterance}  
    & & \multicolumn{4}{p{0.6\linewidth}}{
        \cellcolor[rgb]{0.9,0.9,0.9}{
            \makecell[{{p{\linewidth}}}]{
                \texttt{\tiny{[GM$\rangle$P2]}}
                \texttt{RIFIUTO: \{'A11', 'A60', 'C08', 'B20', 'B33', 'B21', 'B39', 'B96', 'B09'\}} \\
\\ 
\texttt{ARGOMENTAZIONE: \{'La tua proposta continua a includere A11 e A60 che hanno un rapporto valore{-}sforzo problematico. A11 richiede 588 di sforzo mentre A60 ne richiede 255, consumando una porzione significativa del budget. Dovremmo concentrarci su oggetti che offrono maggiore efficienza nella distribuzione delle risorse disponibili.'\}} \\
\\ 
\texttt{PROPOSTA: \{'A82', 'B39', 'B96', 'B09', 'B21', 'B31', 'A21', 'B33', 'C08', 'B20'\}} \\
            }
        }
    }
    & \\ \\

    \theutterance \stepcounter{utterance}  
    & & & \multicolumn{4}{p{0.6\linewidth}}{
        \cellcolor[rgb]{0.9,0.9,0.9}{
            \makecell[{{p{\linewidth}}}]{
                \texttt{\tiny{[GM$\langle$P2]}}
                \texttt{RIFIUTO: \{'A82', 'B39', 'B96', 'B09', 'B21', 'B31', 'A21', 'B33', 'C08', 'B20'\}} \\
\\ 
\texttt{ARGOMENTAZIONE: \{'La tua proposta ha uno sforzo totale di 1539, che lascia ancora 768 unità inutilizzate del nostro budget. Stiamo sprecando opportunità di includere oggetti di alto valore. Inoltre A82 continua ad avere un rapporto valore{-}sforzo molto scarso.'\}} \\
\\ 
\texttt{PROPOSTA: \{'C81', 'A11', 'A60', 'C08', 'B33', 'B21', 'B39', 'B96', 'A21'\}} \\
            }
        }
    }
     \\ \\

    \theutterance \stepcounter{utterance}  
    & & \multicolumn{4}{p{0.6\linewidth}}{
        \cellcolor[rgb]{0.9,0.9,0.9}{
            \makecell[{{p{\linewidth}}}]{
                \texttt{\tiny{[P1$\langle$GM]}}
                \texttt{RIFIUTO: \{'A82', 'B39', 'B96', 'B09', 'B21', 'B31', 'A21', 'B33', 'C08', 'B20'\}} \\
\\ 
\texttt{ARGOMENTAZIONE: \{'La tua proposta ha uno sforzo totale di 1539, che lascia ancora 768 unità inutilizzate del nostro budget. Stiamo sprecando opportunità di includere oggetti di alto valore. Inoltre A82 continua ad avere un rapporto valore{-}sforzo molto scarso.'\}} \\
\\ 
\texttt{PROPOSTA: \{'C81', 'A11', 'A60', 'C08', 'B33', 'B21', 'B39', 'B96', 'A21'\}} \\
            }
        }
    }
    & \\ \\

    \theutterance \stepcounter{utterance}  
    & \multicolumn{4}{p{0.6\linewidth}}{
        \cellcolor[rgb]{0.9,0.9,0.9}{
            \makecell[{{p{\linewidth}}}]{
                \texttt{\tiny{[P1$\rangle$GM]}}
                \texttt{RIFIUTO: \{'C81', 'A11', 'A60', 'C08', 'B33', 'B21', 'B39', 'B96', 'A21'\}} \\
\\ 
\texttt{ARGOMENTAZIONE: \{'Continui a proporre C81 che richiede 543 di sforzo per un valore discutibile, e A11 con 588 di sforzo. Insieme questi due oggetti consumano quasi la metà del nostro budget. Dovremmo trovare un compromesso che includa più oggetti di valore medio{-}alto piuttosto che puntare su pochi oggetti costosi con rendimento incerto.'\}} \\
\\ 
\texttt{PROPOSTA: \{'B93', 'B39', 'B96', 'B09', 'B21', 'B31', 'A21', 'B33', 'A82'\}} \\
            }
        }
    }
    & & \\ \\

    \theutterance \stepcounter{utterance}  
    & & \multicolumn{4}{p{0.6\linewidth}}{
        \cellcolor[rgb]{0.9,0.9,0.9}{
            \makecell[{{p{\linewidth}}}]{
                \texttt{\tiny{[GM$\rangle$P2]}}
                \texttt{RIFIUTO: \{'C81', 'A11', 'A60', 'C08', 'B33', 'B21', 'B39', 'B96', 'A21'\}} \\
\\ 
\texttt{ARGOMENTAZIONE: \{'Continui a proporre C81 che richiede 543 di sforzo per un valore discutibile, e A11 con 588 di sforzo. Insieme questi due oggetti consumano quasi la metà del nostro budget. Dovremmo trovare un compromesso che includa più oggetti di valore medio{-}alto piuttosto che puntare su pochi oggetti costosi con rendimento incerto.'\}} \\
\\ 
\texttt{PROPOSTA: \{'B93', 'B39', 'B96', 'B09', 'B21', 'B31', 'A21', 'B33', 'A82'\}} \\
            }
        }
    }
    & \\ \\

    \theutterance \stepcounter{utterance}  
    & & & \multicolumn{4}{p{0.6\linewidth}}{
        \cellcolor[rgb]{0.9,0.9,0.9}{
            \makecell[{{p{\linewidth}}}]{
                \texttt{\tiny{[GM$\langle$P2]}}
                \texttt{RIFIUTO: \{'B93', 'B39', 'B96', 'B09', 'B21', 'B31', 'A21', 'B33', 'A82'\}} \\
\\ 
\texttt{ARGOMENTAZIONE: \{'Stai ripetendo la stessa combinazione problematica con B93 che consuma 698 unità di sforzo e A82 con rapporto valore{-}sforzo pessimo. Inoltre il totale è solo 1304, sprecando oltre 1000 unità del nostro budget. Dobbiamo essere più strategici nell\textbackslash{}'uso delle risorse disponibili.'\}} \\
\\ 
\texttt{PROPOSTA: \{'C81', 'A60', 'C08', 'B20', 'B31', 'B33', 'B21', 'B39', 'B96', 'A21'\}} \\
            }
        }
    }
     \\ \\

    \theutterance \stepcounter{utterance}  
    & & \multicolumn{4}{p{0.6\linewidth}}{
        \cellcolor[rgb]{0.9,0.9,0.9}{
            \makecell[{{p{\linewidth}}}]{
                \texttt{\tiny{[P1$\langle$GM]}}
                \texttt{RIFIUTO: \{'B93', 'B39', 'B96', 'B09', 'B21', 'B31', 'A21', 'B33', 'A82'\}} \\
\\ 
\texttt{ARGOMENTAZIONE: \{'Stai ripetendo la stessa combinazione problematica con B93 che consuma 698 unità di sforzo e A82 con rapporto valore{-}sforzo pessimo. Inoltre il totale è solo 1304, sprecando oltre 1000 unità del nostro budget. Dobbiamo essere più strategici nell\textbackslash{}'uso delle risorse disponibili.'\}} \\
\\ 
\texttt{PROPOSTA: \{'C81', 'A60', 'C08', 'B20', 'B31', 'B33', 'B21', 'B39', 'B96', 'A21'\}} \\
            }
        }
    }
    & \\ \\

    \theutterance \stepcounter{utterance}  
    & \multicolumn{4}{p{0.6\linewidth}}{
        \cellcolor[rgb]{0.9,0.9,0.9}{
            \makecell[{{p{\linewidth}}}]{
                \texttt{\tiny{[P1$\rangle$GM]}}
                \texttt{RIFIUTO: \{'C81', 'A60', 'C08', 'B20', 'B31', 'B33', 'B21', 'B39', 'B96', 'A21'\}} \\
\\ 
\texttt{ARGOMENTAZIONE: \{'Continuiamo a essere in disaccordo sui grandi oggetti costosi. La tua proposta include ancora C81 che penso sia inefficiente. Proviamo a trovare un compromesso concentrandoci su oggetti di dimensioni medie che potrebbero soddisfare entrambi. Suggerisco di escludere sia i miei oggetti più costosi che i tuoi, e focalizziamoci su una selezione più ampia di oggetti con sforzo moderato.'\}} \\
\\ 
\texttt{PROPOSTA: \{'B39', 'B96', 'B09', 'B21', 'B31', 'B20', 'A21', 'B33', 'C08', 'A60'\}} \\
            }
        }
    }
    & & \\ \\

    \theutterance \stepcounter{utterance}  
    & & \multicolumn{4}{p{0.6\linewidth}}{
        \cellcolor[rgb]{0.9,0.9,0.9}{
            \makecell[{{p{\linewidth}}}]{
                \texttt{\tiny{[GM$\rangle$P2]}}
                \texttt{RIFIUTO: \{'C81', 'A60', 'C08', 'B20', 'B31', 'B33', 'B21', 'B39', 'B96', 'A21'\}} \\
\\ 
\texttt{ARGOMENTAZIONE: \{'Continuiamo a essere in disaccordo sui grandi oggetti costosi. La tua proposta include ancora C81 che penso sia inefficiente. Proviamo a trovare un compromesso concentrandoci su oggetti di dimensioni medie che potrebbero soddisfare entrambi. Suggerisco di escludere sia i miei oggetti più costosi che i tuoi, e focalizziamoci su una selezione più ampia di oggetti con sforzo moderato.'\}} \\
\\ 
\texttt{PROPOSTA: \{'B39', 'B96', 'B09', 'B21', 'B31', 'B20', 'A21', 'B33', 'C08', 'A60'\}} \\
            }
        }
    }
    & \\ \\

    \theutterance \stepcounter{utterance}  
    & & & \multicolumn{4}{p{0.6\linewidth}}{
        \cellcolor[rgb]{0.9,0.9,0.9}{
            \makecell[{{p{\linewidth}}}]{
                \texttt{\tiny{[GM$\langle$P2]}}
                \texttt{ARGOMENTAZIONE: \{'Questa proposta sembra più ragionevole ed elimina gli oggetti più controversi come C81, A11, B93 e A82. Include un buon mix di oggetti a sforzo moderato che potrebbero funzionare per entrambi. Il totale di sforzo è 1539, utilizzando bene il nostro budget senza sprechi eccessivi.'\}} \\
\\ 
\texttt{ACCORDO: \{'B39', 'B96', 'B09', 'B21', 'B31', 'B20', 'A21', 'B33', 'C08', 'A60'\}} \\
            }
        }
    }
     \\ \\

    \theutterance \stepcounter{utterance}  
    & & & \multicolumn{2}{p{0.3\linewidth}}{
        \cellcolor[rgb]{0.9,0.9,0.9}{
            \makecell[{{p{\linewidth}}}]{
                \texttt{\tiny{[GM$|$GM]}}
                \texttt{game successful} \\
            }
        }
    }
    & & \\ \\

    \theutterance \stepcounter{utterance}  
    & & & \multicolumn{2}{p{0.3\linewidth}}{
        \cellcolor[rgb]{0.9,0.9,0.9}{
            \makecell[{{p{\linewidth}}}]{
                \texttt{\tiny{[GM$|$GM]}}
                \texttt{end game} \\
            }
        }
    }
    & & \\ \\

\end{supertabular}
}

\end{document}
